% Autor: Elias Welt
\section{Ausblick und Erweiterungsmöglichkeiten}
Das aktuelle Frontend bietet eine solide Basis, kann jedoch in Zukunft erweitert werden:
\begin{itemize}
    \item \textbf{Automatisierte Tests}: Integration von \texttt{Vitest} für Unit-Tests der Logik und \texttt{Cypress} für End-to-End Tests der Benutzeroberfläche.
    \item \textbf{PWA (Progressive Web App)}: Durch Hinzufügen eines Service Workers und eines Web App Manifests könnte die Anwendung offline-fähig gemacht und auf Mobilgeräten installiert werden.
    \item \textbf{Internationalisierung (i18n)}: Nutzung von Bibliotheken wie \texttt{react-i18next}, um die Oberfläche in mehreren Sprachen anzubieten.
\end{itemize}

\section{Fazit}
Die Frontend-Implementierung von EduShelf demonstriert die Leistungsfähigkeit moderner Webtechnologien. Durch den Einsatz von React und einer komponentenbasierte Architektur entstand eine wartbare, skalierbare und benutzerfreundliche Anwendung. Die Integration komplexer Features wie Real-Time-Editoren, KI-Chat und interaktive Lernmodule in einer Single Page Application zeigt, dass Webanwendungen nativen Programmen in nichts nachstehen müssen.
